\chapter*{General Conclusion}
\addcontentsline{toc}{chapter}{General Conclusion}
\markboth{General Conclusion}{General Conclusion}

The administration of enterprise network security architectures has historically forced a difficult compromise: organizations had to choose between the slow, error-prone safety of manual configuration and the rigid, context-blind speed of static automation. As infrastructure complexity scaled, particularly within the operational realities of integration firms like Kleos Tunisia, this compromise became a critical operational bottleneck. The primary objective of this graduation project was to definitively break this compromise by designing a hybrid orchestration system capable of marrying conversational agility with absolute execution safety.

To achieve this, we departed from the conventional approach of deploying autonomous AI agents, which are inherently susceptible to hallucination and non-deterministic behavior. Instead, we designed and engineered the \textbf{Trust Boundary Architecture}. This architecture physically separated the probabilistic interpretation of administrative intent (driven by the Mistral LLM and ChromaDB RAG pipeline) from the deterministic execution and validation plane (governed by hardcoded Python compliance rules and FortiGate REST API wrappers).

The empirical validation campaign conducted in Chapter 4 demonstrated that this architectural philosophy successfully fulfilled all initial project objectives:
\begin{itemize}
    \item \textbf{Conversational Assistance and Intent Routing:} We successfully bridged the semantic gap between human administrators and device syntax, providing a natural language interface that dynamically classifies and extracts intent without relying on rigid command structures.
    \item \textbf{Contextual Entity Grounding:} By forcing all LLM-extracted entities through a live-state reconciliation engine, we mathematically eliminated the threat of AI hallucination reaching the firewall.
    \item \textbf{Rigid Execution Control:} The implementation of the proactive compliance engine proved that dangerous configurations (e.g., shadowed policies, open \textit{any-to-any} rules) can be deterministically blocked before any network mutation occurs.
    \item \textbf{Human-in-the-Loop and Operational Recoverability:} We guaranteed that the human administrator retains ultimate authority over the infrastructure. The integration of visual delta-diff confirmations, backed by an atomic snapshot and rollback engine, ensured that the system fails safely and is always recoverable.
\end{itemize}

Beyond merely achieving these functional objectives, the system was packaged as a containerized, deployable microservice equipped with a professional Streamlit Security Operations Center (SOC) dashboard and a cryptographic audit chain, elevating it from an academic prototype to a viable enterprise tool.

However, as with any rigorous engineering endeavor, we must acknowledge the limitations of the current implementation. The system currently relies on synchronous API communication, which may introduce latency in WAN-connected environments. Furthermore, while the grounding engine effectively resolves minor semantic ambiguities using Levenshtein distance scoring, deep semantic aliasing remains a challenge that requires future vector-embedding enhancements. 

Looking forward, the Trust Boundary Architecture establishes a highly stable foundation for future innovation. Immediate perspectives for the evolution of this project include extending the orchestration engine to manage fleets of FortiGate appliances via FortiManager integration, rather than a single device. A more advanced horizon involves integrating the agent directly with Security Information and Event Management (SIEM) platforms, allowing the system to autonomously propose—but critically, not execute without human confirmation—quarantine policies in response to active threat intelligence alerts. Additionally, applying formal mathematical verification to the compliance ruleset would provide theoretical guarantees to complement our empirical validation.

In conclusion, this project demonstrated that the integration of artificial intelligence into critical cybersecurity infrastructure does not require the sacrifice of determinism or safety. By strictly constraining probabilistic models within a rigid, deterministic framework, we have established a new paradigm for firewall administration: one that is as agile as human conversation, yet as secure as traditional systems engineering.
